
\section{ОПИСАНИЕ РЕЗУЛЬТАТОВ}

\subsection{Руководство пользователя}

\subsubsection{Роли и интерфейсы}

Система рассчитана на три роли (анализ — 1.1). Каждая роль работает со своим интерфейсом и решает свой класс задач (таблица 4.1); распределение прав доступа ограничивает каждого пользователя точками управления его компетенции.

\begin{xltabular}{\textwidth}{|X|X|X|}
\caption{Соответствие ролей пользователей и интерфейсов системы} \\ \hline
    Роль & Интерфейс & Класс задач \\ \hline
\endfirsthead
\multicolumn{3}{l}{\textit{Продолжение таблицы}} \\ \hline
    Роль & Интерфейс & Класс задач \\ \hline
\endhead
    Контент-менеджер & Веб-интерфейс демо & Выборочная проверка предсказаний, утверждение или отклонение значений атрибутов \\ \hline
    Системный аналитик & Конфигурационные файлы порогов, REST API & Настройка порогов уверенности, мониторинг состояния сервиса \\ \hline
    Инженер машинного обучения & Скрипты обучения, ноутбук с диагностиками & Переобучение классификаторов и байесовских сетей, прогон диагностик, обновление эталона \\ \hline
\end{xltabular}

\subsubsection{Контент-менеджер}

Контент-менеджер работает через веб-интерфейс по адресу \texttt{http://127.0.0.1:8080} после запуска демонстрационного комплекса. Интерфейс состоит из формы ввода партнёр-доступных полей и блока результата с оценкой уверенности и отметкой слоя-источника на каждый атрибут (рис. 4.1).

\begin{figure}
  \includegraphics[width=\textwidth]{fig_4_1_demo_ui.png}
  \caption{Веб-интерфейс демонстрационной системы (схематично): форма ввода партнёр-доступных полей и блок результата с отметкой слоя-источника и оценкой уверенности на каждый атрибут}
  \label{fig:fig_4_1_demo_ui}
\end{figure}

Отметка слоя-источника даёт основание для доверия к предсказанию: слой 1 и слой 4 детерминированы либо объяснимы, для слоя 2 и слоя 3 показателем служит уверенность. Приём введён (см. 1.3) для компенсации ограниченной интерпретируемости слоёв 2/3 на уровне отдельного товара.

\subsubsection{Системный аналитик}

Аналитик настраивает каскад через две точки. Первая — словари порогов \texttt{models/\{category\}\_thresholds.pkl}: контролируют компромисс «точность ↔ покрытие» поатрибутно, изменение не требует переобучения, применяется при следующем старте сервиса (типовой инструмент тонкой настройки при изменении требований к качеству). Вторая — REST API: \texttt{GET /api/categories} (состав атрибутов), \texttt{GET /health} (агрегированное состояние) для интеграции с мониторингом.

\subsubsection{Инженер машинного обучения}

Инженер работает с модулями \texttt{src/} через \texttt{python -m src.<module>} и ноутбуком \texttt{notebooks/00\_thesis\_main.ipynb}. Цикл переобучения устроен так, что выпуск новой версии моделей не требует перекомпиляции каскада: файлы в \texttt{models/} переписываются, ML-сервис при старте загружает обновлённые артефакты. Последовательность команд цикла переобучения:

\begin{verbatim}
source .venv/bin/activate
python -m src.data.label_silver --category pasta_stratified
OMP_NUM_THREADS=1 python -m src.pipeline.ml.train    --category pasta_stratified
OMP_NUM_THREADS=1 python -m src.pipeline.bayes.train --category pasta_stratified
OMP_NUM_THREADS=1 python -m src.eval.run_experiments  --category pasta_stratified
OMP_NUM_THREADS=1 python -m src.eval.run_diagnostics
\end{verbatim}

Цикл запускается еженедельно или ежемесячно (зависит от темпа поступления товаров) и занимает десятки минут на категорию без специализированных ресурсов.

\subsection{Сценарии применения и порядок внедрения}

\subsubsection{Встраивание в бизнес-процесс магазина}

Каскад занимает место промежуточного сервиса в маршруте товара от поставщика до витрины — вторая позиция четырёхзвенной цепочки 1.1, где сосредоточено узкое место. Внедрение не требует перестройки соседних звеньев: ETL партнёрских фидов и PIM работают как раньше, изменяется лишь содержимое «процедуры обогащения». Цепочка обработки:

\begin{verbatim}
Партнёрский фид (CSV/JSON/XML) → ETL магазина (нормализация полей)
  → HTTP /api/enrich — каскад (regex → ML → Bayes → LLM-fallback)
  → Контент-менеджер (выборочная проверка случаев слоя 4)
  → PIM (финальная карточка) → Витрина и поисковый индекс
\end{verbatim}

Основной экономический эффект — перенос ручного труда из режима «закрывать каждое поле» в режим «проверять остаток после каскада». По 3.3.2 каскад покрывает 96,7 \% ячеек без LLM, оставшиеся 3,3 \% уходят в запасной слой; сквозная точность 92,8 \% в конфигурации «каскад + gemini-flash».

\subsubsection{Сценарий холодного старта новой категории}

Для новой категории каскад деградирует до двух слоёв (слой 1 на общих шаблонах + слой 4) — рабочий, но неоптимальный режим: точность слоя 4 на нутриент-производных атрибутах ниже, чем у обученного слоя 2 (4.3.3), стоимость выше. По мере накопления сотен размеченных товаров подключается слой 2; при появлении устойчивых распределений по брендам — слой 3.

Электроника использована как контрольный домен холодного старта (эталон — поля технических характеристик PhoneDB \cite{ref_45} без слабого надзора). Полноценная аудит-оценка по протоколу 3.3 вынесена в раздел «Перспективы развития темы» заключения.

\subsubsection{Сценарий многоязычного каталога}

\texttt{paraphrase-multilingual-MiniLM-L12-v2} обеспечивает единое векторное пространство для пяти доминирующих в OFF европейских языков, что устраняет необходимость отдельных моделей под каждый язык. Устойчивость по языкам подтверждена в 3.3.7.1: разброс между лучшим и худшим сегментом — 2,9 п. п. Каскад применим в международных площадках без удвоения затрат на поддержку.

\subsubsection{Сценарий обогащения нутриентно-производных полей}

Атрибуты, восходящие к числовым нутриентам, требуют точного числового ответа. Одиночный вызов LLM малопригоден из-за высокой частоты ошибок на числах (1.2). В каскаде такие атрибуты закрываются комбинацией слоя 1 (числовые шаблоны) и слоя 2 (классификация дискретных классов после бакетирования нутриентов). USDA Branded Foods используется на этапе обучения для сверки эталона яруса 0 (3.2.5); расширение каскада детерминированным USDA-слоем в режиме вывода — в разделе «Перспективы развития темы» заключения.

\subsubsection{Применимость в смежных доменах}

Архитектура не привязана к пищевому домену: правила слоя 1 переписываются для новой категории, обучение слоёв 2/3 — стандартная процедура. Репликация на электронике (схема атрибутов, серебро из PhoneDB, демонстрация холодного старта; подробности в разделе «Перспективы развития темы» заключения) подтверждает переносимость. Формальная аудит-оценка вынесена в перспективы развития.

\subsection{Что позволяет использование разработки (характеристика достигнутых результатов)}

Использование системы даёт пять измеримых эффектов в формате «повышено / ускорено / сокращено» со ссылками на 3.

\subsubsection{Повышено: качество автоматического обогащения каталога}

Сквозная точность на 4350 ячейках без пересечения брендов — 92,8 \% в конфигурации «каскад + gemini-flash» (3.3.2; \texttt{grand\_acc\_summary\_after\_fix.parquet}). 95 \%-й интервал Вильсона: [92,0; 93,6] для «каскад + gemini-flash» и [89,6; 91,4] для «только каскад» (90,5 \%). Поатрибутный разрез по категориям приведён в таблице 4.2 (источник: \texttt{headline\_v3e\_after\_fix.parquet}).

\begin{xltabular}{\textwidth}{|X|X|X|X|}
\caption{Сквозная точность каскада по категориям (модель gemini-flash, 4350 ячеек без пересечения брендов)} \\ \hline
    Категория & Только каскад & Каскад + gemini-flash & Прирост \\ \hline
\endfirsthead
\multicolumn{4}{l}{\textit{Продолжение таблицы}} \\ \hline
    Категория & Только каскад & Каскад + gemini-flash & Прирост \\ \hline
\endhead
    Паста & 93,3 \% & 95,8 \% & +2,5 п. п. \\ \hline
    Шоколад & 87,4 \% & 90,1 \% & +2,7 п. п. \\ \hline
    Сыры & 90,4 \% & 92,2 \% & +1,8 п. п. \\ \hline
    Взвешенное среднее & 90,5 \% & 92,8 \% & +2,3 п. п. \\ \hline
\end{xltabular}

\textit{Примечание. Источник: \texttt{headline\_v3e\_after\_fix.parquet}.}

Сыры и шоколад получают наибольший выигрыш от запасного слоя на оставшихся ≈3 \% сложных ячеек. На пасте, где правила слоя 1 покрывают значительную долю, прирост от слоя 4 минимален.

Сравнение с альтернативами: прямой Claude Sonnet 4.5 даёт 83,8 \% (на 9,0 п. п. ниже системы); прямой gpt-oss-120b — 69,3 \% (на 23,0 п. п. ниже каскада с этой же моделью в роли слоя 4 — 92,3 \%). Превосходство объясняется наличием специализированных слоёв: Слой 2 обучен на нутриент-производных атрибутах (Nutri-Score, класс белка, класс жирности), которые LLM не вычисляет точно из текста без доступа к числам.

\subsubsection{Сокращено: стоимость обработки в 720 раз относительно эталонной LLM-модели}

Стоимость обработки одного товара через каскад + gemini-flash составляет ≈0,14 \% стоимости прямого Sonnet 4.5 (сокращение в 720 раз при превышении точности на 9,0 п. п.; матрица — 3.3.2). Экономия — за счёт обработки полей в порядке возрастания стоимости: Слой 1 (бесплатно) → Слой 2 (миллисекунды на CPU) → Слой 3 (локальный вывод) → Слой 4 только при отказе предыдущих. По 3.3 LLM вызывается на 3,3 \% ячеек, а не на 100 \%. Стоимость пяти каскадных конфигураций относительно опорной «все ячейки — Sonnet 4.5» приведена в таблице 4.3 (источник: \texttt{grand\_acc\_summary\_after\_fix.parquet}).

\begin{xltabular}{\textwidth}{|X|X|X|}
\caption{Стоимость обработки одного товара относительно опорной конфигурации all-sonnet} \\ \hline
    Конфигурация & Стоимость относительно опорной & Снижение стоимости \\ \hline
\endfirsthead
\multicolumn{3}{l}{\textit{Продолжение таблицы}} \\ \hline
    Конфигурация & Стоимость относительно опорной & Снижение стоимости \\ \hline
\endhead
    Каскад + gemini-flash & 0,0014 & в 720 раз \\ \hline
    Каскад + gpt-4o & 0,024 & в 42 раза \\ \hline
    Каскад + Claude Sonnet 4.5 & 0,033 & в 30 раз \\ \hline
    Каскад + gpt-oss-120b (открытая модель) & 0,0010 & в 1007 раз \\ \hline
\end{xltabular}

\textit{Примечание. Источник: \texttt{grand\_acc\_summary\_after\_fix.parquet}.}

При сценарии 100 000 товаров/месяц экономия составляет порядки величины относительно прямого вызова LLM.

\subsubsection{Обеспечена: возможность развёртывания на собственных мощностях с открытой моделью}

Система поддерживает gpt-oss-120b в роли слоя 4 — полностью локальное развёртывание без отправки партнёрских данных вовне. Сводно (3.3.2): сквозная точность каскад + gpt-oss-120b — 92,3 \%; прямой gpt-oss-120b — 69,3 \%; прирост от каскадной архитектуры — +23,0 п. п. (слой 2 закрывает атрибуты, на которых малая модель проигрывает); стоимость сценария ≈в 30 раз ниже прямого использования за счёт сокращения LLM-вызовов до 3,3 \% ячеек. При потере 0,5 п. п. точности относительно gemini-flash обеспечивается независимость от внешних API — критично для корпоративных сценариев с требованиями локального хранения данных.

\subsubsection{Сокращён: ручной труд оператора каталога}

В исходном процессе (1.1) оператор проверяет каждый атрибут — 100 \% ручной работы. Система переводит труд в режим выборочной верификации: 96,7 \% ячеек закрываются дешёвыми слоями автоматически; 3,3 \% уходят в слой 4 и получают предсказание (часть требует верификации, обычно — помеченные байесовский валидатором); совокупная доля ручной верификации — 3–7 \% в зависимости от порога валидатора (3.3.4). Ручная нагрузка сокращается в 10–20 раз; контроль качества сохраняется через отметку слоя-источника и доверительные индикаторы.

\subsubsection{Повышена: устойчивость каталога к смещению по брендам}

Тестирование без пересечения брендов (3.2.4) подтверждает обобщение на бренды, не встречавшиеся в обучении: на 4350 ячейках средняя точность 92,8 \% с интервалом Вильсона [92,0; 93,6]. Дополнительно k-NN-абляция (3.3.7.3) показывает: точность не падает на удалённых от обучения товарах, а монотонно растёт от Q1 (82,7 \%) до Q4 (91,6 \%).

\subsubsection{Резюмирующая таблица «повышено / ускорено / сокращено»}

Сводный эффект разработки на семи показателях приведён в таблице 4.4 (источники: \texttt{grand\_acc\_summary\_after\_fix.parquet}, \texttt{headline\_v3e\_after\_fix.parquet}).

\begin{xltabular}{\textwidth}{|X|X|X|X|}
\caption{Сводный эффект разработки на ключевых показателях системы (повышено / сокращено / обеспечено)} \\ \hline
    Показатель & До разработки & После & Эффект \\ \hline
\endfirsthead
\multicolumn{4}{l}{\textit{Продолжение таблицы}} \\ \hline
    Показатель & До разработки & После & Эффект \\ \hline
\endhead
    Сквозная точность системы & 83,8 \% (опорная конфигурация all-sonnet) & 92,8 \% (каскад + gemini-flash) & повышена на 9,0 п. п. \\ \hline
    Относительная стоимость обработки & 1,000 (опорная конфигурация all-sonnet) & 0,0014 (каскад + gemini-flash) & сокращена в 720 раз \\ \hline
    Доля обращений к LLM & 100 \% & 3,3 \% & сокращена в 30 раз \\ \hline
    Ручная верификация оператором & ~100 \% ячеек & ~5–10 \% ячеек & сокращена в 10–20 раз \\ \hline
    Точность на открытой модели & 69,3 \% (одиночный вызов gpt-oss-120b) & 92,3 \% (каскад + gpt-oss-120b) & повышена на 23,0 п. п. \\ \hline
    Поддержка многоязычности & отдельные модели на каждый язык либо предварительный перевод & 5 рабочих языков в едином векторном пространстве SBERT & обеспечена без дополнительных моделей \\ \hline
    Независимость от поставщика LLM & один интегрированный провайдер & поддержка пяти моделей с сопоставимой точностью & обеспечена взаимозаменяемость \\ \hline
\end{xltabular}

\textit{Примечание. Источники: \texttt{grand\_acc\_summary\_after\_fix.parquet}, \texttt{headline\_v3e\_after\_fix.parquet}.}

Значения столбцов «До» и «После» получены на тестовой выборке без пересечения брендов главы 3 и подтверждены артефактами \texttt{grand\_acc\_summary\_after\_fix.parquet} и \texttt{headline\_v3e\_after\_fix.parquet}.

\subsection{Выводы по главе 4}

\begin{enumerate}[arabic]
    \item Система допускает три сценария по выбору слоя 4: бюджетный коммерческий (gemini-flash) — наилучший баланс точности и стоимости; премиум-коммерческий (Sonnet 4.5, GPT-4o) — близкая точность при большей стоимости; локальное развёртывание с gpt-oss-120b — независимость от внешних API при потере 0,5 п. п. точности.
    \item Целевая аудитория — три роли: контент-менеджер (выборочная верификация через веб-интерфейс), системный аналитик (пороги через конфигурации и REST API), инженер ML (переобучение по мере поступления данных).
    \item Совокупный эффект: сокращение стоимости в 720 раз относительно прямого вызова LLM и ручной нагрузки в 10–20 раз при точности 92,8 \% на выборке без пересечения брендов с интервалом Вильсона [92,0; 93,6].
\end{enumerate}
